\section{Coefficient budgets, coordinate restrictions, and fixed channels}
\label{p1:arith:supplement}

The first two tools below do not depend on the critical-window scaling.
They control, respectively, the small-prime coefficients of selected
values and the repetition of coordinate forms in a binary relation space.
We then give the extension estimates for any fixed rational channel.
In that part, the channel and every finite family of channels are fixed
before \(N\) tends to infinity.

\subsection{A deterministic budget for the small-prime part}

For a positive integer \(n\), write
\[
 u_B(n)=\frac{\sf(n)}{K_B(n)}
       =\prod_{\substack{\ell\leq B\ {\rm prime}\\v_\ell(n)\ {\rm odd}}}\ell.
\]

\begin{lemma}[A small-prime coefficient budget]
\label{p1:arith:smallpart-budget}
Let \(V\) be any interval of \(B\geq2\) consecutive positive integers.
Then
\[
 \sum_{n\in V}\log u_B(n)\leq B\log B+O(B),
\]
with an absolute constant independent of the position of \(V\).
For every fixed \(A>0\), at most
\[
 \bigl(A^{-1}+O_A(1/\log B)\bigr)B
\]
of its values satisfy \(u_B(n)>B^A\).
\end{lemma}

\begin{proof}
Dropping the parity condition gives
\[
 \sum_{n\in V}\log u_B(n)
 \leq\sum_{\ell\leq B}
       \left(\left\lfloor\frac B\ell\right\rfloor+1\right)\log\ell
 \leq\log(B!)+\sum_{\ell\leq B}\log\ell.
\]
Here the sums are over primes.  Legendre's formula gives the second
inequality, and the elementary Chebyshev bound makes the last sum
\(O(B)\).  Finally \(\log(B!)\leq B\log B\).
Each value counted in the last assertion contributes more than
\(A\log B\) to the nonnegative sum.
\end{proof}

\subsection{Regularising coordinate forms by independent restrictions}

\begin{lemma}[Coordinate regularisation]
\label{p1:arith:coordinate-regularisation}
Let \(R\subset\F_2^V\) have dimension \(r\), put \(n=|V|\), and let
\(\ell_v\in R^*\) be the coordinate forms.  Suppose that \(s\) distinct
nonzero forms occur among them.  For integers \(K,d\geq1\), there is
\(D\subset V\) such that
\[
 R'=\{x\in R:x_v=0\text{ for every }v\in D\},\qquad
 \operatorname{codim}_R R'=|D|
 \leq\frac{n}{K+1}+\frac{s-r}{d+1},
\]
and the restricted coordinate forms have the following properties:
\begin{enumerate}
 \item every nonzero form has multiplicity at most \(K\);
 \item every nonzero form \(\ell\) has at most \(d\) unordered
 decompositions \(\ell=a+b\), where \(a,b\) are distinct nonzero
 coordinate forms present on \(R'\).
\end{enumerate}
Coordinates whose restricted forms vanish may then be removed without
any additional loss of dimension.
\end{lemma}

\begin{proof}
The coordinate forms span \(R^*\), and their restrictions continue to
span the dual of every subspace encountered below.
At an intermediate stage let \(S\) be the set of distinct nonzero forms
and let \(r_0\) be their rank.  Setting a present nonzero form \(\ell\)
equal to zero reduces the dimension by one and replaces the forms by
their images modulo \(\langle\ell\rangle\).
The form \(\ell\) disappears.  All other collisions are exactly the
disjoint unordered pairs \(\{a,a+\ell\}\subset S\).
If there are \(t(\ell)\) such pairs, then
\[
 |S_{\rm new}|=|S|-1-t(\ell),\qquad
 (|S|-r_0)_{\rm new}=|S|-r_0-t(\ell).
\]
Thus the nonnegative potential \(|S|-r_0\) never increases.

If a form has multiplicity greater than \(K\), set it equal to zero.
This permanently kills at least \(K+1\) previously nonzero coordinates.
There are at most \(n/(K+1)\) such steps, since these sets of killed
coordinates are disjoint.
When the multiplicity condition holds but a form has more than \(d\)
decompositions, set that form equal to zero instead.
Such a step decreases the potential by at least \(d+1\), so there are
at most \((s-r)/(d+1)\) of these steps.
Each step decreases the dimension, and the procedure therefore ends
with both required properties.

At every step select one coordinate carrying the form just killed.
Its restriction was nonzero before that step, so its original coordinate
form is independent of the previously selected forms.
The selected coordinates constitute \(D\), and consequently
\(\operatorname{codim}_R R'=|D|\).
\end{proof}

\begin{corollary}[Constraints meeting three blocks]
\label{p1:arith:ternary-degree}
Suppose \(V\) is partitioned into three blocks and each row of a family
of binary constraints annihilating \(R\) has at most one occurrence in
each block.  Apply Lemma~\ref{p1:arith:coordinate-regularisation} and
delete the coordinates whose forms on \(R'\) are zero.
Every remaining row has weight zero, two, or three.
A row of weight two joins equal nonzero forms.
A row of weight three has three distinct nonzero forms \(a,b,c\) with
\(a+b+c=0\).

The connected components formed by the rows of weight two have at most
\(K\) vertices.  After repeated rows with the same support are identified,
the number of ternary rows incident to any vertex is at most \(2dK^2\).
\end{corollary}

\begin{proof}
The sum of the forms on a row is zero, since that row annihilates \(R'\).
A row of weight one is therefore impossible.  At weight two its forms
are equal.  At weight three a repeated form would force the third form
to be zero, which has already been excluded.
Binary components stay inside a class of equal forms, of size at most
\(K\).  At a vertex of form \(\ell\), a ternary row chooses one of at
most \(d\) decompositions \(\ell=a+b\), then the other two coordinates
from classes of size at most \(K\).  Allowing their two assignments to
the remaining blocks gives the asserted \(2dK^2\) bound.
\end{proof}

\begin{remark}[Order of restrictions]
If some coordinates must first be killed for an arithmetic reason, do
this before applying the regularisation lemma.  A further restriction
after regularisation can identify previously distinct forms and destroy
both bounds.  The lemma is an existence statement about a subspace;
its proof does not enumerate possible relation spaces or histories of
restrictions.
\end{remark}

\subsection{A finite rank bound for every fixed channel}

Fix coprime integers \(a,b\) with
\[
 b\geq3,\qquad b/2<a<b.
\]
Let \(\mathcal G_{a:b}\) be the symmetric family of ordered separated
pairs in this channel with at least two exact units and with a private
prime coordinate on every distinct normalized value in
\(bV_x\cup aV_y\), above a fixed threshold \(H_{a,b}=C_{a,b}B\).
Choose the constant large enough for the fixed normalized widths.
For \(\alpha=(x,y)\) put
\[
 q_\alpha=\P(J_x=J_y=1),\qquad
 M_{a:b}=\sum_{\alpha\in\mathcal G_{a:b}}q_\alpha.
\]
Pairs of probability zero remain in the family and are omitted from
conditional expressions.  The private-prime condition makes the
relative kernel exactly the canonical even-selection kernel.

\begin{lemma}[Mass of a fixed channel]
\label{p1:arith:fixed-channel-mass}
With \(S_{a,b}(L)\) as defined before
\eqref{p1:arith:channelidentity}, and
\(w_{a,b}=2(2a-b)/(ab)\), one has
\[
 M_{a:b}=p^2Nw_{a,b}S_{a,b}(L)
          +O_{a,b,C,\varepsilon}(N^{-1+\varepsilon}).
\]
In particular \(M_{a:b}\asymp_{a,b,C}L N^{-1+1/b}\).
\end{lemma}

\begin{proof}
In either orientation the canonical parameter satisfies
\(N/a\leq t<2N/b\), up to fixed multiples of \(B\) at the endpoints.
The two orientations therefore supply \(Nw_{a,b}\) translations per
choice of the within-window placements.
The diameter and sign calculation giving
\eqref{p1:arith:channelidentity} applies for this fixed \(a,b\).

Removing the normalized defective pairs costs
\(N^{-3/2+1/b+o(1)}\): there are \(N^{1/2+o(1)}\) such hosts,
their nullity is at most \(L/b+o(L)\), and their weights include \(p^2\).
Endpoint errors are smaller.
On the retained dependent pairs the signed kernel contribution is
\(q_\alpha-p^2\).  Adding back \(p^2\) costs \(O_{a,b}(p^2NB)\),
because each start has \(O_{a,b}(B)\) possible partners.
This proves the displayed formula.
Its leading term is positive and of order \(L N^{-1+1/b}\), by the
finite identity, uniformly in the critical band.
\end{proof}

Define the finite, explicitly computable number
\begin{equation}\label{p1:arith:gamma-fixed-channel}
 \gamma_{a:b}=
 \max\left(
 \left\{\frac1b\right\}\cup
 \left\{\frac{(c,a)+(c,b)-1}{c}:
 \begin{array}{l}
 1\leq c\leq b(a+b-1),\ d\geq1,\ (c,d)=1,\\
 b/2<c/d<2a,\quad c\nmid a,\quad c\nmid b
 \end{array}\right\}\right).
\end{equation}
All variables in the second set are integers.

\begin{lemma}[Additional rank in a fixed channel]
\label{p1:arith:fixed-channel-rank}
One has \(\gamma_{a:b}\leq2/3\).
Apart from \(B^{O_{a,b}(1)}\) source pairs at the dyadic endpoints,
every aligned target \(z\ne x,y\) compatible with a source
\((x,y)\in\mathcal G_{a:b}\) satisfies
\[
 \delta_{x,y}(z):=
 L-\dim\frac{S_x+S_y+S_z}{S_x+S_y}
 \leq\gamma_{a:b}L+o_{a,b,C}(L).
\]
This definition includes any internal rank deficiency of the new window.
The bound remains true when the prime signs of \(a,b\) are included in
the old information.
\end{lemma}

\begin{proof}
Write \(x=at-i_0\), \(y=bt-j_0\), choosing the lower orientation.
As in Lemma~\ref{p1:arith:palmrank}, the normalized old values lie in
a fixed-width interval about \(abt\).  An aligned target has the form
\(z=(ct+h)/d\), with \((c,d)=1\), numerator, denominator and offsets
of polynomial size in \(B\).
The dyadic positions require
\(b/2\leq c/d\leq2a\), up to \(B^{O(1)}/N\).
Separation of rational fractions excludes strict violations of this
interval once \(N\) is large.
Either equality forces \(at=N+B^{O(1)}\) or \(bt=2N+B^{O(1)}\);
these give the stated exceptional sources.

For the interior fractions, a new offset \(k\) matched to a true old
value satisfies one of
\[
 h+dk\equiv0\pmod{c/(c,a)},\qquad
 h+dk\equiv0\pmod{c/(c,b)}.
\]
Their intersection has modulus \(c\).
The union contains at most
\[
 B\,\frac{(c,a)+(c,b)-1}{c}+O(1)
\]
new positions.  The same residual argument as in
Lemma~\ref{p1:arith:palmrank} leaves only \(o(L)\) further dimensions.
Matching to a value in an enclosing interval but absent from the true
old set contributes no new freedom.

Suppose \(c\nmid a\) and \(c\nmid b\), and put \(g=(c,a)\), \(h_1=(c,b)\).
These numbers are coprime and \(gh_1\mid c\).
If one equals one, the other is a proper divisor of \(c\), giving
density at most \(1/2\).
Otherwise \(g,h_1\geq2\) and
\[
 \frac{g+h_1-1}{c}
 \leq\frac1g+\frac1{h_1}-\frac1{gh_1}\leq\frac23.
\]
The last maximum occurs at the smallest coprime pair \(\{2,3\}\).
When \(c>b(a+b-1)\), the density is less than \(1/b\).
All other nonphysical fractions are included in
\eqref{p1:arith:gamma-fixed-channel}.

If \(c\mid a\), a proper divisor has \(c\leq a/2<b/2\), which is
impossible here; thus \(c=a,d=1\).
Similarly \(c\mid b\) leaves only \(c=b,d=1\).
These are the physical fractions reproducing one old family.
Compatibility gives distance at least \(L\) from that source, so the
two true windows share at most one value.
Matches to the other old family have steps \(a,b\) in intervals of
width \(B\), and number at most \(B/b+O(1)\).
This proves the rank bound in these cases as well.
\end{proof}

\begin{proposition}[Palm mean for a fixed channel]
\label{p1:arith:fixed-channel-palm}
Put
\[
 \eta_{a:b}=\min\{17/24,1-\gamma_{a:b}\}\geq1/3.
\]
Then, with the sum over \(z\in I_N\setminus\{x,y\}\),
\[
 R_{a:b}:=
 \frac1{M_{a:b}}\sum_{\alpha=(x,y)\in\mathcal G_{a:b}}q_\alpha
 \sum_z\bigl(\P(J_z=1\mid J_x=J_y=1)-p\bigr)_+
 \ll_{a,b,C,\varepsilon}N^{-\eta_{a:b}+\varepsilon}.
\]
Moreover
\[
 \frac1{M_{a:b}}\sum_{\alpha=(x,y)\in\mathcal G_{a:b}}q_\alpha
 \E[Z_N-2\mid J_x=J_y=1]
 =\lambda+
 O_{a,b,C,\varepsilon}
 \bigl(N^{-\min(\eta_{a:b},1/2)+\varepsilon}\bigr).
\]
These bounds are not asserted to be optimal, or uniform in growing
channel height.
\end{proposition}

\begin{proof}
The source probability is \(O_{a,b,C}(p^2 2^{L/b})\).
For the compressed old interval choose a fixed enclosing radius \(DB\).
Use this same radius in the alignment test and the complementary
aligned-target sector; the large-prime cutoff is chosen accordingly.
The nonaligned asymmetric estimate
in Lemma~\ref{p1:arith:asymmetric}, with the fixed larger widths and
\(B^{O_{a,b}(1)}\) patterns, bounds the total positive excess by
\[
 p^3 2^{L/b}N^{31/24+\varepsilon'}.
\]
Division by \(M_{a:b}\asymp NLp^2 2^{L/b}\) gives
\(N^{-17/24+\varepsilon}\).

Each source has \(B^{O_{a,b}(1)}\) aligned targets.
Lemma~\ref{p1:arith:fixed-channel-rank} bounds their positive conditional
excess by \(pN^{\gamma_{a:b}+o(1)}\) per target.
The exceptional endpoint sources have total weight
\(B^{O(1)}p^2 2^{L/b}\).
Even bounding every conditional probability by one, their aligned
contribution after division by \(M_{a:b}\) is \(N^{-1+o(1)}\).
This proves the first assertion, including all new windows.

For the mean, Proposition~\ref{p1:arith:fixed-branching} with two
exposed source windows leaves at most \(N^{1/2+o(1)}\) exceptional
targets for each source.
Outside them the new relative vector is uniform and independent of
the old values, so its conditional probability is exactly \(p\).
The total negative deviation is therefore at most \(pN^{1/2+o(1)}\).
Combining it with the positive estimate and the removal of the two
sources, which costs \(2p\), proves the assertion.
\end{proof}

\subsection{The ternary mass generated by a finite family of pairs}

\begin{corollary}[Triples containing a good pair]
\label{p1:arith:finite-family-triples}
Let \(\mathcal F\) be a nonempty finite fixed family of admissible
channels, and set
\[
 M_{\mathcal F}=\sum_{a:b\in\mathcal F}M_{a:b},\qquad
 \nu_{\mathcal F}=
 \min\{1/2,\min_{a:b\in\mathcal F}\eta_{a:b}\}.
\]
The sum of \(q(\{x,y,z\})\) over all ordered triples of distinct starts
containing at least one good pair from this family is
\[
 3\lambda M_{\mathcal F}
 +O_{\mathcal F,C,\varepsilon}
    (M_{\mathcal F}N^{-\nu_{\mathcal F}+\varepsilon}).
\]
In particular, for fixed \(Q\geq5\), the triples containing a good
pair of height \(5\leq b\leq Q\) have total mass
\(O_{Q,C,\varepsilon}(N^{-4/5+\varepsilon})\).
\end{corollary}

\begin{proof}
Choose which two positions of an ordered triple form the designated
pair.  For each of the three choices,
Proposition~\ref{p1:arith:fixed-channel-palm} gives
\(\lambda M_{\mathcal F}
 +O(M_{\mathcal F}N^{-\nu_{\mathcal F}+\varepsilon})\).
The pair orientations are already included in \(M_{\mathcal F}\).

This designated-pair sum counts a triple once for each good pair that
it contains.  If the third window is independent of all old values
of the designated pair, it cannot form a dependent good pair with
either source.
Every overcount therefore has its third start in the exceptional set
of Proposition~\ref{p1:arith:fixed-branching}.
For each designated source, the conditional mass of this set is at most
\[
 pN^{1/2+o(1)}
 +\sum_z\bigl(\P(J_z=1\mid J_x=J_y=1)-p\bigr)_+.
\]
Its weighted average is
\(O_{\mathcal F,C,\varepsilon}(N^{-\nu_{\mathcal F}+\varepsilon})\).
This bounds the entire overcount, since a triple has at most three pairs.
The finite channel families are disjoint for large \(N\), by the separation
of distinct reduced slopes.
Finally, Lemma~\ref{p1:arith:fixed-channel-mass} gives
\(M_{\mathcal F}\ll_{Q,C}L N^{-4/5}\) for heights between five and \(Q\).
\end{proof}

The last corollary concerns the part of the third moment containing a
good pair.  It makes no assertion about the primitive remainder or about
a family whose heights grow with \(N\).

